% ============================================================================
%  Modelo de relatório — Iscte-IUL
%  Compilar SEMPRE este ficheiro (os restantes têm "% !TEX root = ../main.tex")
% ============================================================================
\documentclass[a4paper,11pt,oneside,openany]{report}

% !TEX root = ../main.tex
% ============================================================================
%  Informação do relatório — é praticamente só isto que muda entre relatórios
% ============================================================================

\newcommand{\universidade}{Iscte -- Instituto Universitário de Lisboa}
\newcommand{\escola}{Instituto Superior de Ciências do Trabalho e da Empresa}
\newcommand{\curso}{Mestrado em Cibersegurança e Resiliência}
\newcommand{\uc}{Nome da Unidade Curricular}

\newcommand{\titulo}{Título do Relatório}
\newcommand{\subtitulo}{Trabalho Prático n.º 1}
\newcommand{\titulocurto}{Título Curto}   % aparece no cabeçalho das páginas

% Uma linha por autor:  Nome & \email{...} \\
\newcommand{\autores}{%
  Nome Apelido  & \email{nome_apelido@iscte-iul.pt} \\
  Outro Nome    & \email{outro_nome@iscte-iul.pt} \\
}
\newcommand{\grupo}{Grupo 00}             % deixar vazio {} se não aplicável
\newcommand{\docente}{Nome do Docente}

\newcommand{\anoletivo}{Ano Letivo 2026/2027 \textbar{} 1.º Semestre}
\newcommand{\datarelatorio}{\today}

% Logótipo da capa: branco com fundo transparente (vai sobre o azul)
\newcommand{\ficheirologo}{imagens/logo-iscte.png}

% Cor principal (títulos, legendas, capa) — azul Iscte
\newcommand{\corhex}{0E29C2}
        % título, UC, autores, docente, datas  <-- editar
% !TEX root = ../main.tex
% ============================================================================
%  Pacotes
% ============================================================================

% --- Língua e tipografia ---------------------------------------------------
\usepackage[T1]{fontenc}
\usepackage[utf8]{inputenc}
\usepackage[english, portuguese]{babel}   % a última é a língua principal
\usepackage{csquotes}
\usepackage[sfdefault]{plex-sans}       % IBM Plex Sans (texto)
\usepackage[scale=0.95]{plex-mono}      % IBM Plex Mono (código, emails)
% Alternativas (trocar a linha do plex-sans):
%   \usepackage[sfdefault]{FiraSans}        % mais técnico/compacto
%   \usepackage[default]{sourcesanspro}     % mais neutro
%   \usepackage{plex-serif}                 % texto com serifa, mais clássico
\usepackage{microtype}

% --- Página -----------------------------------------------------------------
\usepackage[a4paper, margin=2.5cm, headheight=15pt]{geometry}
\usepackage{setspace}
\usepackage{indentfirst}
\usepackage{fancyhdr}
\usepackage{titlesec}
\usepackage[titles]{tocloft}
\usepackage[nottoc, notbib]{tocbibind}

% --- Figuras e tabelas ------------------------------------------------------
\usepackage{xcolor}
\usepackage{tikz}               % capa
\usepackage{graphicx}
\graphicspath{{imagens/}}       % \includegraphics{topologia} procura em imagens/
\usepackage{float}              % permite [H]
\usepackage{caption}
\usepackage{subcaption}
\usepackage{booktabs}
\usepackage{tabularx}
\usepackage{longtable}
\usepackage{multirow}
\usepackage{makecell}
\usepackage{pdfpages}           % \includepdf para anexar PDFs (ex.: enunciado)

% --- Matemática, unidades, listas ------------------------------------------
\usepackage{amsmath, amssymb}
\usepackage{siunitx}
\sisetup{output-decimal-marker={,}}
\usepackage{enumitem}

% --- Código -----------------------------------------------------------------
\usepackage{listings}

% --- Referências ------------------------------------------------------------
% style=ieee (numérico, comum em engenharia). Para APA: style=apa
\usepackage[backend=biber, style=ieee]{biblatex}

% --- Ligações (hyperref perto do fim; glossaries depois dele) ---------------
\usepackage{hyperref}
\usepackage[nameinlink]{cleveref}
\usepackage[acronym, toc, nonumberlist, nopostdot, nomain]{glossaries}
\makenoidxglossaries            % não precisa de makeglossaries / perl
     % pacotes
% !TEX root = ../main.tex
% ============================================================================
%  Estilo
% ============================================================================

\definecolor{corprincipal}{HTML}{\corhex}
\definecolor{fundocodigo}{gray}{0.96}
\colorlet{azulclaro}{corprincipal!35!white}

\newcommand{\email}[1]{\href{mailto:#1}{\nolinkurl{#1}}}

\onehalfspacing
\setlength{\parindent}{1.5em}
\setlist{itemsep=0.2em}

\hypersetup{
  colorlinks = true,
  linkcolor  = black,
  citecolor  = corprincipal,
  urlcolor   = corprincipal,
  pdftitle   = {\titulo},
}

% --- Nomes em português -----------------------------------------------------
\addto\captionsportuguese{%
  \renewcommand{\contentsname}{Índice}%
  \renewcommand{\listfigurename}{Índice de Figuras}%
  \renewcommand{\listtablename}{Índice de Tabelas}%
  \renewcommand{\appendixname}{Anexo}%
}
\renewcommand{\lstlistingname}{Código}
\renewcommand{\lstlistlistingname}{Índice de Código}

\newcommand{\crefrangeconjunction}{ a\nobreakspace}
\newcommand{\crefpairconjunction}{ e\nobreakspace}
\newcommand{\crefmiddleconjunction}{, }
\newcommand{\creflastconjunction}{ e\nobreakspace}
\crefname{figure}{figura}{figuras}        \Crefname{figure}{Figura}{Figuras}
\crefname{table}{tabela}{tabelas}          \Crefname{table}{Tabela}{Tabelas}
\crefname{equation}{equação}{equações}     \Crefname{equation}{Equação}{Equações}
\crefname{chapter}{capítulo}{capítulos}    \Crefname{chapter}{Capítulo}{Capítulos}
\crefname{section}{secção}{secções}        \Crefname{section}{Secção}{Secções}
\crefname{subsection}{secção}{secções}     \Crefname{subsection}{Secção}{Secções}
\crefname{appendix}{anexo}{anexos}         \Crefname{appendix}{Anexo}{Anexos}
\crefname{listing}{código}{códigos}        \Crefname{listing}{Código}{Códigos}
\crefname{lstlisting}{código}{códigos}     \Crefname{lstlisting}{Código}{Códigos}

% --- Títulos ----------------------------------------------------------------
\titleformat{\chapter}[hang]
  {\normalfont\huge\bfseries\color{corprincipal}}{\thechapter}{0.6em}{}
\titlespacing*{\chapter}{0pt}{0pt}{1.2em}
\titleformat{\section}
  {\normalfont\Large\bfseries\color{corprincipal}}{\thesection}{0.8em}{}
\titleformat{\subsection}
  {\normalfont\large\bfseries\color{corprincipal}}{\thesubsection}{0.8em}{}
\titleformat{\subsubsection}
  {\normalfont\normalsize\bfseries\color{corprincipal}}{\thesubsubsection}{0.8em}{}

% --- Legendas ---------------------------------------------------------------
\captionsetup{labelfont={bf, color=corprincipal}, font=small, labelsep=period}

% --- Índice -----------------------------------------------------------------
\renewcommand{\cftchapleader}{\cftdotfill{\cftdotsep}}
\setlength{\cftbeforechapskip}{0.4em}

% --- Cabeçalhos e rodapés ---------------------------------------------------
\fancypagestyle{plain}{%            páginas de início de capítulo / pré-textuais
  \fancyhf{}%
  \fancyfoot[C]{\thepage}%
  \renewcommand{\headrulewidth}{0pt}%
}
\fancypagestyle{corpo}{%
  \fancyhf{}%
  \fancyhead[L]{\small\color{corprincipal}\nouppercase{\leftmark}}%
  \fancyhead[R]{\small\color{corprincipal}\titulocurto}%
  \fancyfoot[C]{\thepage}%
  \renewcommand{\headrulewidth}{0.4pt}%
}
\renewcommand{\chaptermark}[1]{\markboth{\thechapter.\ #1}{}}

% --- Código -----------------------------------------------------------------
\lstset{
  backgroundcolor   = \color{fundocodigo},
  basicstyle        = \ttfamily\footnotesize,
  keywordstyle      = \color{corprincipal}\bfseries,
  commentstyle      = \color{gray}\itshape,
  stringstyle       = \color{teal!70!black},
  numbers           = left,
  numberstyle       = \tiny\color{gray},
  frame             = single,
  rulecolor         = \color{gray!40},
  breaklines        = true,
  tabsize           = 2,
  showstringspaces  = false,
  captionpos        = b,
  % listings não suporta UTF-8: sem isto, acentos dentro de código dão erro
  literate =
    {á}{{\'a}}1 {à}{{\`a}}1 {â}{{\^a}}1 {ã}{{\~a}}1
    {é}{{\'e}}1 {ê}{{\^e}}1 {í}{{\'i}}1
    {ó}{{\'o}}1 {ô}{{\^o}}1 {õ}{{\~o}}1 {ú}{{\'u}}1 {ü}{{\"u}}1
    {Á}{{\'A}}1 {À}{{\`A}}1 {Â}{{\^A}}1 {Ã}{{\~A}}1
    {É}{{\'E}}1 {Ê}{{\^E}}1 {Í}{{\'I}}1
    {Ó}{{\'O}}1 {Ô}{{\^O}}1 {Õ}{{\~O}}1 {Ú}{{\'U}}1
    {ç}{{\c c}}1 {Ç}{{\c C}}1 {º}{{\textordmasculine}}1 {ª}{{\textordfeminine}}1
    {€}{{\texteuro}}1 {°}{{\textdegree}}1,
}

% --- Acrónimos --------------------------------------------------------------
\newcommand{\printacronimos}{%
  \printnoidxglossary[type=\acronymtype, title={Lista de Acrónimos}, toctitle={Lista de Acrónimos}]}

% --- Transições de numeração ------------------------------------------------
\newcommand{\iniciarcorpo}{%
  \clearpage
  \pagenumbering{arabic}%
  \pagestyle{corpo}%
}

\newcommand{\iniciaranexos}{%
  \appendix
  \titleformat{\chapter}[hang]
    {\normalfont\huge\bfseries\color{corprincipal}}{\appendixname~\thechapter}{0.6em}{}%
  \renewcommand{\chaptermark}[1]{\markboth{\appendixname~\thechapter.\ ##1}{}}%
  \addtocontents{toc}{%
    \protect\renewcommand{\protect\cftchappresnum}{\appendixname~}%
    \protect\setlength{\protect\cftchapnumwidth}{5.2em}}%
  \renewcommand{\thepage}{\thechapter-\arabic{page}}%
}

% Usar \anexo{Título} em vez de \chapter nos anexos: reinicia a página (A-1, B-1, ...)
\newcommand{\anexo}[1]{\chapter{#1}\setcounter{page}{1}}
      % aspeto: cores, títulos, cabeçalhos, código, anexos
% !TEX root = ../main.tex
% ============================================================================
%  Acrónimos — usar no texto com \gls{chave}
%  (1.ª vez: "Border Gateway Protocol (BGP)"; seguintes: "BGP")
%  Só aparecem na lista os que forem usados no texto.
% ============================================================================
\setacronymstyle{long-short}

\newacronym{bgp}{BGP}{Border Gateway Protocol}
\newacronym{as}{AS}{Autonomous System}
\newacronym{ospf}{OSPF}{Open Shortest Path First}
   % lista de acrónimos                   <-- editar
\addbibresource{referencias.bib}

\begin{document}

% ---------------------------------------------------------------------------
%  Elementos pré-textuais — numeração romana (i, ii, iii, ...)
% ---------------------------------------------------------------------------
% !TEX root = ../main.tex
% Capa: bloco azul Iscte (logótipo + título) e zona branca (autores, docente, data)
\begin{titlepage}
\begin{tikzpicture}[remember picture, overlay]
  \coordinate (azulfim) at ([yshift=-0.58\paperheight]current page.north west);

  % bloco azul
  \fill[corprincipal] (current page.north west) rectangle (azulfim -| current page.east);

  % grelha de pontos a desvanecer (canto superior direito)
  \foreach \x in {0,...,11} {
    \foreach \y in {0,...,13} {
      \pgfmathsetmacro{\op}{max(0, 0.30 - 0.025*\x - 0.012*\y)}
      \fill[white, opacity=\op]
        ([xshift={-1.6cm - \x*0.55cm}, yshift={-1.6cm - \y*0.55cm}]current page.north east)
        circle (1.1pt);
    }
  }

  % logótipo
  \node[anchor=north west, inner sep=0]
    at ([xshift=2.5cm, yshift=-2.4cm]current page.north west) {%
    \IfFileExists{\ficheirologo}
      {\includegraphics[width=5cm]{\ficheirologo}}
      {\color{white}\Huge\bfseries iscte}};

  % UC, título, subtítulo
  \node[anchor=south west, inner sep=0, align=left, text width=16cm]
    at ([xshift=2.5cm, yshift=1.8cm]azulfim) {%
    {\color{azulclaro}\ttfamily\small > \curso\ /\ \uc\par}
    \vspace{0.6cm}
    {\color{white}\fontsize{30}{36}\selectfont\bfseries \titulo\par}
    \vspace{0.4cm}
    {\color{white}\Large \subtitulo\par}};

  % autores
  \node[anchor=north west, inner sep=0, align=left, text width=10cm]
    at ([xshift=2.5cm, yshift=-1.5cm]azulfim) {%
    {\color{corprincipal}\footnotesize\bfseries\textls[150]{AUTORES}%
      \ifx\grupo\empty\else\ \ \textperiodcentered\ \ \textls[150]{\MakeUppercase{\grupo}}\fi\par}
    \vspace{0.3cm}
    \begin{tabular}{@{}l@{\hspace{1.2em}}l@{}}
      \autores
    \end{tabular}};

  % docente e data
  \node[anchor=north east, inner sep=0, align=right, text width=5.5cm]
    at ([xshift=-2.5cm, yshift=-1.5cm]azulfim -| current page.east) {%
    {\color{corprincipal}\footnotesize\bfseries\textls[150]{DOCENTE}\par}
    \vspace{0.3cm}
    \docente\par
    \vspace{0.8cm}
    {\color{corprincipal}\footnotesize\bfseries\textls[150]{DATA}\par}
    \vspace{0.3cm}
    \datarelatorio\par};

  % rodapé
  \draw[corprincipal, line width=1.5pt]
    ([xshift=2.5cm, yshift=2.6cm]current page.south west) --
    ([xshift=-2.5cm, yshift=2.6cm]current page.south east);
  \node[anchor=north west, inner sep=0, font=\small, text=gray]
    at ([xshift=2.5cm, yshift=2.3cm]current page.south west) {\escola};
  \node[anchor=north east, inner sep=0, font=\small, text=gray]
    at ([xshift=-2.5cm, yshift=2.3cm]current page.south east) {\anoletivo};
\end{tikzpicture}
\null
\end{titlepage}


\pagenumbering{roman}
\pagestyle{plain}

% !TEX root = ../main.tex
\chapter*{Resumo}
\addcontentsline{toc}{chapter}{Resumo}

Breve descrição do trabalho (150--250 palavras): contexto, objetivo, o que foi
feito, principais resultados e conclusão.

\vspace{1em}
\noindent\textbf{Palavras-chave:} palavra1, palavra2, palavra3

% !TEX root = ../main.tex
\chapter*{Abstract}
\addcontentsline{toc}{chapter}{Abstract}

\begin{otherlanguage}{english}
Short description of the work in English.

\vspace{1em}
\noindent\textbf{Keywords:} keyword1, keyword2, keyword3
\end{otherlanguage}
   % comentar se não for pedido abstract em inglês

\tableofcontents
\listoffigures                % comentar se não houver figuras
\listoftables                 % comentar se não houver tabelas
\lstlistoflistings            % comentar se não houver código
\printacronimos               % comentar se não houver acrónimos

% ---------------------------------------------------------------------------
%  Corpo do relatório — numeração árabe (1, 2, 3, ...)
% ---------------------------------------------------------------------------
\iniciarcorpo

% !TEX root = ../main.tex
\chapter{Introdução}
\label{cap:introducao}

\section{Contexto}

Enquadramento do trabalho. Neste relatório estuda-se o \gls{bgp}, o protocolo
usado para trocar rotas entre cada \gls{as} da Internet~\cite{rfc4271}.

\section{Objetivos}

\begin{itemize}
  \item Objetivo 1;
  \item Objetivo 2;
  \item Objetivo 3.
\end{itemize}

\section{Estrutura do relatório}

O \cref{cap:desenvolvimento} descreve o trabalho realizado, o
\cref{cap:conclusao} apresenta as conclusões e o \cref{anx:configuracoes}
contém as configurações completas.

% !TEX root = ../main.tex
\chapter{Desenvolvimento}
\label{cap:desenvolvimento}

Este capítulo serve também de \emph{cábula}: mostra como fazer cada elemento.
Apagar ou substituir à vontade.

\section{Figuras}

As imagens ficam em \texttt{imagens/} e não é preciso escrever a pasta nem a
extensão. Referência automática: \cref{fig:exemplo}.

\begin{figure}[H]
  \centering
  \includegraphics[width=0.6\linewidth]{example-image-a} % ex.: {topologia}
  \caption{Legenda da figura}
  \label{fig:exemplo}
\end{figure}

Duas figuras lado a lado (\cref{fig:lado-a,fig:lado-b}):

\begin{figure}[H]
  \centering
  \begin{subfigure}{0.45\linewidth}
    \centering
    \includegraphics[width=\linewidth]{example-image-b}
    \caption{Antes}
    \label{fig:lado-a}
  \end{subfigure}
  \hfill
  \begin{subfigure}{0.45\linewidth}
    \centering
    \includegraphics[width=\linewidth]{example-image-c}
    \caption{Depois}
    \label{fig:lado-b}
  \end{subfigure}
  \caption{Comparação antes/depois}
\end{figure}

\section{Tabelas}

\begin{table}[H]
  \centering
  \caption{Endereçamento IP}   % nas tabelas a legenda vai em cima
  \label{tab:enderecos}
  \begin{tabular}{llll}
    \toprule
    \textbf{Router} & \textbf{Interface} & \textbf{IP privado} & \textbf{IP público} \\
    \midrule
    R1 & Lo0  & 10.1.1.1/32 & --- \\
       & g1/0 & 10.1.2.1/30 & --- \\
    R2 & Lo1  & ---         & 47.73.239.22/32 \\
    \bottomrule
  \end{tabular}
\end{table}

Para tabelas que ocupam várias páginas usar \texttt{longtable}.

\section{Código}

\begin{lstlisting}[language=Java, caption={Exemplo de Web Service}, label={lst:exemplo}]
@WebService
public class Servico {
    @WebMethod
    public String ola(String nome) {
        return "Olá, " + nome + "! Já funcionam acentos.";
    }
}
\end{lstlisting}

Código num ficheiro externo:
\verb|\lstinputlisting[language=Python]{codigo/script.py}|.
Configuração de um router, sem linguagem: \verb|\begin{lstlisting}...\end{lstlisting}|.

\section{Equações}

\begin{equation}
  \label{eq:exemplo}
  T = \frac{L}{R} = \frac{\SI{1500}{\byte}}{\SI{100}{\mega\bit\per\second}}
\end{equation}

Ver \cref{eq:exemplo}. Os acrónimos já usados aparecem curtos: \gls{bgp}, \gls{as}.
Livro de referência:~\cite{kurose2021}.

% !TEX root = ../main.tex
\chapter{Conclusão}
\label{cap:conclusao}

Síntese do que foi feito, resultados obtidos face aos objetivos, dificuldades
encontradas e trabalho futuro.


\printbibliography[heading=bibintoc, title={Referências}]

% ---------------------------------------------------------------------------
%  Anexos — capítulos A, B, C; páginas A-1, A-2, B-1, ...
% ---------------------------------------------------------------------------
\iniciaranexos

% !TEX root = ../main.tex
\anexo{Configurações completas}
\label{anx:configuracoes}

Nos anexos as figuras e tabelas ficam numeradas A.1, A.2, \ldots{} e as páginas
A-1, A-2, \ldots

\begin{lstlisting}[caption={Configuração do R1}]
interface Loopback0
 ip address 10.1.1.1 255.255.255.255
!
router ospf 1
 network 10.0.0.0 0.255.255.255 area 0
\end{lstlisting}

\begin{figure}[H]
  \centering
  \includegraphics[width=0.5\linewidth]{example-image}
  \caption{Figura num anexo}
\end{figure}

% !TEX root = ../main.tex
\anexo{Enunciado}
\label{anx:enunciado}

Para anexar um PDF inteiro (ex.: o enunciado), colocá-lo em \texttt{anexos/} e usar:

\verb|\includepdf[pages=-, pagecommand={}]{anexos/enunciado.pdf}|


\end{document}
